\chapter{Generics}
\label{ch:generics}

You met \code{Vector<int>} back in Chapter~7 and the bound \code{<T: Shape>} in
Chapter~11. Both rely on \textbf{generics}: the ability to write code that works
for \emph{any} type, with the type filled in later. Generics let you write a
function or a data structure once and reuse it for integers, strings, your own
structs anything without copying and pasting. This chapter brings the idea
into full focus.

\section{The problem generics solve}

Suppose you want a function that returns the larger of two values. For integers:

\begin{salam}
func max_int(a : int, b : int) : int:
    if a > b: ret a end
    ret b
end
\end{salam}

For floating-point numbers you would write \emph{the same logic again} with
\code{f64} in place of \code{int}, and again for any other comparable type. That
is wasteful and error-prone a bug fixed in one copy lingers in the others.
Generics let you write the logic a single time.

\section{Generic functions}

A \textbf{type parameter} in angle brackets after the function name turns a
function generic. The parameter conventionally \code{T} stands for ``some
type, chosen by the caller,'' and you use it like any other type within the
function:

\begin{salam}
func Max<T>(a : T, b : T) : T:
    if a > b: ret a end
    ret b
end

func main:
    println Max(10, 20)        // T is int
    println Max(2.5, 1.5)      // T is f64
end
\end{salam}

\begin{progout}
20
2.5
\end{progout}

One \code{Max} now serves every type. The caller writes \code{Max(10, 20)} with
no angle brackets; Salam looks at the arguments and \emph{infers} that \code{T} is
\code{int} for the first call and \code{f64} for the second. The two parameters
\code{a} and \code{b} share the type \code{T}, so they must be the \emph{same}
type at each call you cannot write \code{Max(10, 2.5)}, because \code{T} cannot
be both \code{int} and \code{f64} at once.

\section{Generic structs}

Data structures benefit even more from generics. A ``box that holds one value''
should not care whether that value is an integer or a string. Declare a type
parameter after the struct's name and use it for fields:

\begin{salam}
struct Box<T>:
    value : T
end

func main:
    mut a : Box<int> = Box { value = 42 }
    mut b : Box<str> = Box { value = "hi" }
    println a.value
    println b.value
end
\end{salam}

\begin{progout}
42
hi
\end{progout}

\code{Box<int>} is a box of integers; \code{Box<str>} is a box of strings. You
spell out the concrete type in the variable's declaration \code{Box<int>} ---
and from then on the field \code{value} has that type. This is precisely how
\code{Vector<T>} works: it is a generic struct holding a growable buffer of
\code{T}.

\begin{note}
At the \emph{call site} of a generic function, you usually omit the type ---
\code{Max(10, 20)} because Salam infers it from the arguments. But when you
declare a variable of a generic \emph{struct} type, you write the concrete type
explicitly: \code{Box<int>}, \code{Vector<str>}. The struct's type parameter
cannot always be inferred from an empty \code{Box \{\}}, so you state it.
\end{note}

\section{A generic container: a simple stack}

Let us build something real: a generic stack, a last-in-first-out store, on top of
a fixed array. Because it is generic, the same code gives you a stack of
integers, a stack of strings, or a stack of anything:

\begin{salam}
struct Stack<T>:
    items : T[16] = [0,0,0,0,0,0,0,0,0,0,0,0,0,0,0,0]
    count : int = 0
    func push(x : T):
        this.items[this.count] = x
        this.count = this.count + 1
    end
    func pop() : T:
        this.count = this.count - 1
        ret this.items[this.count]
    end
    func is_empty() : bool:
        ret this.count == 0
    end
end

func main:
    mut s : Stack<int> = Stack {}
    s.push(3)
    s.push(7)
    s.push(1)
    println s.pop()
    println s.pop()
    println s.is_empty()
end
\end{salam}

\begin{progout}
1
7
false
\end{progout}

The methods \code{push} and \code{pop} are written entirely in terms of \code{T}.
Declare \code{Stack<int>} and you get an integer stack; declare \code{Stack<str>}
and the very same definition gives you a string stack. This is how the standard
library's collections are built, and how you build your own.

\begin{note}
The \code{items} field carries a default (sixteen zeros) so that \code{Stack
\{\}} is a complete value with every field initialized the same rule from
Chapter~9 that every struct field needs a default if you want to construct the
struct with empty braces. For a real, growable container you would build on
\code{Vector<T>} (Chapter~7) instead of a fixed array; this small fixed-size
stack keeps the focus on the generics.
\end{note}

\section{Bounded generics revisited}

A bare \code{<T>} accepts \emph{any} type, but then you can only do to a \code{T}
what every type supports. The moment you call a method on a \code{T} ---
\code{s.area()}, say the compiler must know that \code{T} actually \emph{has}
that method. That is what an interface \textbf{bound} (Chapter~11) provides:
\code{<T: Shape>} means ``any \code{T}, as long as it conforms to \code{Shape},''
which unlocks the interface's methods inside the function.

\begin{salam}
interface Shape:
    func area() : f64
end

func total_area<T: Shape>(a : T, b : T) : f64:
    ret a.area() + b.area()
end
\end{salam}

So there is a spectrum: \code{<T>} for code that treats values opaquely (store,
move, return them), and \code{<T: SomeInterface>} for code that needs to
\emph{do} something specific with them. You reach for the bound exactly when you
need to call the interface's methods.

\section{How generics work: monomorphization}

\begin{deep}[In Depth: one template, many concrete copies]
Salam implements generics by \emph{monomorphization}, the same mechanism behind
interface bounds. The generic definition \code{Max<T>} or \code{Stack<T>} ---
is a \emph{template}; the compiler never emits it directly. Instead, each time you
use it with a concrete type, the compiler stamps out a specialised copy with the
type filled in: \code{Max} becomes a real \code{Max\_i32} and a real
\code{Max\_f64} (the compiler names the copies with the explicit type names);
\code{Stack<int>} becomes a concrete C struct with an \code{int[16]} field. Those
copies are ordinary, fully type-checked code, as fast
as anything hand-written there is no boxing, no run-time type tag, no
indirection. The trade-off is compiled size: one copy exists for each distinct
type you instantiate. In practice that is a small price for code that is both
reusable and fast.

To make this possible, the compiler needs to know an element's size at compile
time. The built-in \code{sizeof(T)} gives it: the standard library's
\code{Vector<T>}, for instance, uses \code{sizeof(T)} to allocate exactly enough
bytes for each element when it grows its buffer. You rarely call \code{sizeof}
yourself, but it is the quiet machinery that lets one generic container hold
elements of any size.
\end{deep}

\section{When to make something generic}

Generics are powerful, but not everything needs them. A good rule:

\begin{itemize}[leftmargin=1.4em]
  \item Make a \emph{container} generic almost always a stack, queue, list, or
        tree is naturally indifferent to what it stores.
  \item Make a \emph{function} generic when you find yourself about to copy it for
        a second type. The first copy is fine; the second is the signal.
  \item Do \emph{not} reach for generics when a single concrete type is all you
        will ever need the extra abstraction earns its keep only when there is
        real reuse.
\end{itemize}

\section{Chapter summary}

\begin{itemize}[leftmargin=1.4em]
  \item Generics write code once for \emph{any} type, using a type parameter
        \code{<T>}.
  \item A generic function infers \code{T} from its arguments at the call site;
        parameters sharing \code{T} must be the same type.
  \item A generic struct (\code{Box<T>}, \code{Stack<T>}, \code{Vector<T>}) is
        declared with its concrete type, e.g. \code{Stack<int>}.
  \item A bare \code{<T>} treats values opaquely; an interface bound \code{<T:
        Iface>} lets you call the interface's methods on a \code{T}.
  \item Generics compile by \emph{monomorphization} into specialised, fast copies;
        \code{sizeof(T)} is the machinery that lets a container size its elements.
\end{itemize}

\section{Exercises}

\exercise Write a generic function \code{Min<T>(a T, b T) T} mirroring \code{Max},
and test it on integers and floats.

\exercise Write a generic function \code{first\_of<T>(a T, b T) T} that simply
returns its first argument, and explain why it needs no bound.

\exercise Define a generic \code{struct Pair<T>} holding two values of type
\code{T} and a method \code{swap()} that exchanges them. Test it with a pair of
integers.

\exercise Extend the generic \code{Stack} with a \code{peek()} method that returns
the top element without removing it. Test it with a \code{Stack<str>}.

\exercise Explain in a comment why \code{Max(10, 2.5)} fails to compile, in terms
of the type parameter \code{T}.
